% #################################################################
% JGME AI Teaching Rounds Manuscript Template
% Modeled on:
%   "Supervising Resident AI Use Without Losing the Learning"
%   Journal of Graduate Medical Education, February 2026
%
% #################################################################

\documentclass[12pt]{article}

\usepackage[margin=1in]{geometry}
\usepackage{setspace}
\usepackage{enumitem}
\usepackage{booktabs}
\usepackage{longtable}
\usepackage{tabularx}
\usepackage{graphicx}
\usepackage{hyperref}
\usepackage{url}
\usepackage{microtype}
\usepackage{fancyhdr}
\usepackage{titlesec}
\usepackage{xcolor}

\hypersetup{
  hidelinks,
  pdfauthor={AUTHOR NAME},
  pdftitle={AI TEACHING ROUNDS MANUSCRIPT TITLE}
}

\doublespacing

\pagestyle{fancy}
\fancyhf{}
\fancyfoot[C]{\thepage}

\setlength{\parindent}{0.5in}
\setlength{\parskip}{0pt}
\setlist[itemize]{leftmargin=*}
\setlist[enumerate]{leftmargin=*}

% ------------------------------------------------------------------
% USER VARIABLES
% ------------------------------------------------------------------

\newcommand{\ArticleTitle}{FULL AI TEACHING ROUNDS TITLE}
\newcommand{\AuthorName}{AUTHOR NAME, DEGREE(S)}
\newcommand{\AuthorAffiliation}{DEPARTMENT, INSTITUTION, CITY, STATE/PROVINCE, COUNTRY}
\newcommand{\CorrespondingEmail}{EMAIL}
\newcommand{\CaseTitle}{The Case of [SHORT, MEMORABLE CASE TITLE]}

\begin{document}

% ==================================================================
% TITLE / AUTHOR
% ==================================================================

\begin{center}
	{\Large\bfseries \ArticleTitle\par}

	\vspace{0.8em}

	\AuthorName
\end{center}

\vspace{1em}

% ==================================================================
% OPENING CASE
% ==================================================================

\section*{\CaseTitle}

 % PURPOSE:
 % Open with a concrete, believable GME scenario involving a resident,
 % faculty member, program, or training environment.
 %
 % MODEL:
 % The published article begins with an overnight ED case involving a
 % second-year resident, an ambient AI scribe, polished documentation,
 % and a gap between documentation quality and observable reasoning.
 %
 % WRITING TARGET:
 % - 2-5 short paragraphs.
 % - Use a recognizable clinical/educational tension.
 % - End with a practical supervisory or educational question.
 % - Avoid making the case merely decorative: it should introduce the
 %   central problem that the rest of the article resolves.

 [Describe a realistic residency/fellowship scenario.]

 [Show where AI enters the trainee's work.]

 [Show the educational tension: efficiency vs learning, assistance vs
  independence, polish vs reasoning, etc.]

 [End with the central question:
  What should the attending/program do?]

 % ==================================================================
 % CONNECTION TO THE SERIES
 % ==================================================================

 % The model article explicitly connects back to the inaugural AI Teaching
 % Rounds principles before stating the focus of the current installment.

 [OPTIONAL SERIES BRIDGE.]

In prior discussions of AI in graduate medical education, three principles
are especially useful: normalize disclosure of AI use, prioritize observable
reasoning over polished products, and begin with small, testable steps that
can be revised as programs learn.

In this installment of AI Teaching Rounds, we apply those principles to
	[YOUR SPECIFIC PROBLEM].

% ==================================================================
% WHY THIS MATTERS NOW
% ==================================================================

\section*{Why This Matters Now}

 % PURPOSE:
 % Explain why the problem is immediate rather than hypothetical.
 %
 % The model article does three things:
 % 1. Establishes that AI use is rapidly becoming routine.
 % 2. Contrasts health-system efficiency goals with educational goals.
 % 3. Re-examines assumptions about which clinical tasks are actually
 %    educationally valuable.
 %
 % Suggested 3-part structure:
 %
 % Paragraph 1: adoption / urgency.
 % Paragraph 2: educational risk or tension.
 % Paragraph 3: what this forces programs to define explicitly.

 [Describe the current state of AI adoption relevant to your topic.]

 [Explain why the issue creates a GME-specific educational challenge.]

 [Describe the tension between operational efficiency and development of
  trainee knowledge, judgment, reasoning, communication, or professional
  identity.]

 [State what programs/faculty now need to decide.]

% ==================================================================
% DEFINE THE PROTECTED HUMAN WORK
% ==================================================================

\section*{What Should Remain Human-Only (for Now)}

 % PURPOSE:
 % Identify the parts of the task that trainees should perform before AI
 % assistance because the struggle itself is educationally valuable.
 %
 % The model article protects:
 % - first-pass differential diagnosis;
 % - initial problem representation;
 % - assessment and plan;
 % while allowing AI as a later check or augmentation.
 %
 % Adapt this to your own domain.

 [Identify 2-5 cognitive, communicative, procedural, reflective, or
  professional tasks that should remain trainee-generated first.]

For each protected task, explain:

\begin{itemize}
	\item what competency the task develops;
	\item what would be lost if AI performed it first;
	\item when AI may appropriately enter afterward;
	\item whether the rule should vary by learner stage.
\end{itemize}

% ==================================================================
% PRACTICAL SUPERVISION / TEACHING FRAMEWORK
% ==================================================================

\section*{Supervising AI Use in [YOUR CONTEXT]: The [FRAMEWORK-NAME] Framework}

% PURPOSE:
% Give readers something memorable and usable in real time.
%
% The model article introduces DEFT-AI and then explains each component
% through concrete supervisory questions.
%
% If adapting an existing framework:
% - cite it clearly;
% - state that you are applying/adapting it;
% - do not imply authorship of the framework.
%
% If proposing your own framework:
% - make each step operational;
% - keep the number of steps small;
% - connect each step to the educational risk identified earlier.

\subsection*{Step 1: [NAME]}

[Explain the first action.]

\textit{Example supervisor prompt:}
``[QUESTION THE SUPERVISOR CAN ACTUALLY ASK.]''

\subsection*{Step 2: [NAME]}

[Explain how faculty should inspect inputs, evidence, assumptions, or
context.]

\textit{Example supervisor prompt:}
``[QUESTION.]''

\subsection*{Step 3: [NAME]}

[Explain how the supervisor elicits the trainee's own reasoning.]

\textit{Example supervisor prompt:}
``[QUESTION.]''

\subsection*{Step 4: [NAME]}

[Explain verification, skepticism, checking, or adjudication.]

\textit{Example supervisor prompt:}
``[QUESTION.]''

\subsection*{Step 5: [NAME]}

[Explain how the supervisor gives advice about future appropriate AI use.]

\textit{Example supervisor prompt:}
``[QUESTION.]''

% ==================================================================
% BOX 1
% ==================================================================

\section*{Box 1. [FRAMEWORK NAME] in Under 2 Minutes for [TARGET USER]}

\begin{description}[style=nextline,leftmargin=0pt]
	\item[\textbf{1. [STEP NAME]}]
	      ``[Short question or action.]''

	\item[\textbf{2. [STEP NAME]}]
	      ``[Short question or action.]''

	\item[\textbf{3. [STEP NAME]}]
	      ``[Short question or action.]''

	\item[\textbf{4. [STEP NAME]}]
	      ``[Short question or action.]''

	\item[\textbf{5. [STEP NAME]}]
	      ``[Short question or action.]''
\end{description}

% ==================================================================
% TASK DISSECTION
% ==================================================================

\section*{Task Dissection: What Exactly Is AI Offloading, and Do We Want That?}

 % PURPOSE:
 % Move from individual supervision to program-level analysis.
 %
 % The model article asks:
 % - What is this task teaching?
 % - What exactly would AI offload?
 % - Is that offloading helpful or harmful at this learner stage?
 % - What new learning opportunity might AI create?
 % - How will supervisors still see the trainee's thinking?

 [Explain why real-time supervision alone is insufficient.]

 [Introduce a task-by-task method for deciding whether AI use should be
  protected, augmented, or offloaded.]

% ==================================================================
% BOX 2
% ==================================================================

\section*{Box 2. Task Dissection Framework: 5 Questions}

\begin{enumerate}
	\item What is this task teaching?
	\item What exactly would the AI offload?
	\item For this trainee at this stage, is that offloading helpful or harmful?
	\item What new learning opportunity might AI create?
	\item How will supervisors still see the trainee's thinking?
\end{enumerate}

% ==================================================================
% APPLIED EXAMPLES
% ==================================================================

% The published article then applies the framework to three common tasks.
% Use 2-4 examples from your own educational domain.

\section*{[Task Example 1]}

 [What the task teaches.]

 [What AI offloads.]

 [What should remain trainee-generated.]

 [How AI can be used afterward.]

 [What the supervisor should observe or ask.]

\section*{[Task Example 2]}

 [Repeat the same reasoning.]

\section*{[Task Example 3]}

 [Repeat the same reasoning.]

% ==================================================================
% CROSS-SPECIALTY / CROSS-PROGRAM TABLE
% ==================================================================

\section*{Table. Examples of ``Human-First'' Tasks Across Residency Programs}

\begin{longtable}{p{0.20\textwidth} p{0.36\textwidth} p{0.36\textwidth}}
	\toprule
	\textbf{Residency / Setting}
	 &
	\textbf{Example ``Human-First'' Task}
	 &
	\textbf{How AI Might Be Used After the Resident Attempts}
	\\
	\midrule
	\endfirsthead

	\toprule
	\textbf{Residency / Setting}
	 &
	\textbf{Example ``Human-First'' Task}
	 &
	\textbf{How AI Might Be Used After the Resident Attempts}
	\\
	\midrule
	\endhead

	[Specialty 1]
	 &
	[Human-first cognitive task]
	 &
	[Permitted AI augmentation after trainee attempt]
	\\

	[Specialty 2]
	 &
	[Human-first cognitive task]
	 &
	[Permitted AI augmentation after trainee attempt]
	\\

	[Specialty 3]
	 &
	[Human-first cognitive task]
	 &
	[Permitted AI augmentation after trainee attempt]
	\\

	[Specialty 4]
	 &
	[Human-first cognitive task]
	 &
	[Permitted AI augmentation after trainee attempt]
	\\

	\bottomrule
\end{longtable}

% ==================================================================
% SMALL EXPERIMENTS
% ==================================================================

\section*{Box 3. Try This: Small Experiments for [YOUR TOPIC]}

% The model article provides three very small, immediately actionable
% interventions. Follow that pattern.

\textbf{Experiment 1: [SHORT LABEL]}

[One sentence describing the behavior.]

[One sentence explaining what to observe.]

\vspace{0.6em}

\textbf{Experiment 2: [SHORT LABEL]}

[One sentence describing the behavior.]

[One sentence explaining what to observe.]

\vspace{0.6em}

\textbf{Experiment 3: [SHORT LABEL]}

[One sentence describing the behavior.]

[One sentence explaining what to observe.]

% ==================================================================
% USEFUL RESOURCES
% ==================================================================

\section*{Box 4. Useful Resources}

% The published article does more than list citations: each selected source
% has a short annotation explaining why it is practically useful.

\textbf{Author AA, Author BB. Title. Journal. Year;volume(issue):pages.
	doi:...}

[1-3 sentences explaining why this resource is useful to GME educators.]

\vspace{0.8em}

\textbf{Organization. Guideline/resource title. URL}

[1-3 sentences explaining what practical question this resource helps answer.]

\vspace{0.8em}

\textbf{Author AA, et al. Title. Journal. Year;volume(issue):pages.
	doi:...}

[1-3 sentences explaining its practical relevance.]

% ==================================================================
% CONCLUSION
% ==================================================================

\section*{Conclusions}

 % The model conclusion is short and returns to the main tension introduced
 % in the opening case.
 %
 % Recommended pattern:
 % 1. State what does NOT need to be abandoned.
 % 2. State what DOES need to become explicit.
 % 3. Reaffirm trainee reasoning/human judgment.
 % 4. End with a constructive rather than alarmist view of AI.

 [2-4 concise paragraphs.]

 [Return explicitly to the educational problem introduced in the opening.]

 [State what faculty/programs should protect.]

 [State what AI can appropriately offload or augment.]

 [End with the core educational principle.]

% ==================================================================
% REFERENCES
% ==================================================================

\section*{References}

% Follow JGME's numbered biomedical reference style.
% Number references consecutively in order of first citation.

\begin{enumerate}
	\item Author AA, Author BB. Article title. \textit{Journal Name}.
	      YEAR;VOLUME(ISSUE):PAGES. doi:DOI

	\item Organization. Guideline/resource title. Accessed MONTH DAY, YEAR.
	      \url{https://example.org}

	\item Author AA, Author BB, Author CC, et al. Article title.
	      \textit{Journal Name}. YEAR;VOLUME(ISSUE):PAGES. doi:DOI
\end{enumerate}

% ==================================================================
% AUTHOR INFORMATION
% ==================================================================

\vspace{1em}

\noindent
\AuthorName\ is [ROLE/TITLE], \AuthorAffiliation.

\vspace{0.5em}

\noindent
\textbf{Corresponding author:}
\AuthorName,
\AuthorAffiliation,
\CorrespondingEmail

% ==================================================================
% WRITING CHECKLIST
% ==================================================================

\clearpage
\section*{AI Teaching Rounds Structural Checklist}

\begin{itemize}[label=$\square$]
	\item Title states the educational problem rather than merely naming AI.
	\item Opening case contains a real GME decision or supervisory tension.
	\item Case ends with a question that the article subsequently answers.
	\item ``Why This Matters Now'' establishes urgency and GME relevance.
	\item Human-first/protected learning tasks are explicitly identified.
	\item Article distinguishes AI augmentation from AI substitution.
	\item Learner stage is considered where relevant.
	\item A concise, memorable framework is provided or appropriately adapted.
	\item Framework contains concrete questions/actions faculty can use.
	\item At least one box converts concepts into a rapid-use tool.
	\item Task-dissection section explains what AI actually offloads.
	\item 2-4 applied examples demonstrate the framework.
	\item Cross-specialty table is included if it materially helps readers.
	\item Small, low-stakes experiments are offered.
	\item Practical resources are annotated, not merely listed.
	\item Privacy/PHI and institution-approved tool use are addressed if relevant.
	\item Verification and skepticism are explicitly taught.
	\item Human reasoning remains observable to supervisors.
	\item Conclusion returns to the opening educational tension.
	\item Tone is practical and non-hyperbolic.
	\item References use numbered JGME-style citations.
\end{itemize}

\end{document}
